\iffalse

La teoría de Morse es una poderosa herramienta dentro de la topología diferencial. Permite usar técnicas básicas de análisis para describir una variedad. En este trabajo desarrollamos la teoría de Morse utilizando como fuente principal el libro ``Morse Theory'' de John Milnor \cite{milnor1963}. Seguido de una breve introducción a las variedades diferenciables, presentamos el objeto central de la teoría, las funciones de Morse. Definimos el concepto de punto crítico no degenerado y demostramos el lema de Morse. Posteriormente, enunciamos y demostramos los principales resultados que permiten relacionar el tipo de homotopía de una variedad diferenciable con el de un CW-complejo. Incluimos un capítulo dedicado a la teoría de homolgía para concluir con la demostración de las desigualdades de Morse, que nos permiten acotar en gran medida los números de Betti de una variedad.

%%%%%%%%%%%%%%%%%%%%%%%%%%%%

La teoría de Morse permite extraer de una variedad $M$ importante información topológica mediante el estudio de los puntos críticos de una función diferenciable $f:M\to \mathbb{R}$. En la presente memoria desarrollamos los fundamentos de la teoría de Morse con el libro ``Morse Theory'' de John Milnor \cite{milnor1963} como fuente principal.

Diremos que una función de una variedad diferenciable en los reales es de Morse si todos sus puntos críticos son no degenerados. Demostraremos el lema de Morse y, gracias a él, que los puntos críticos no degenerados son aislados. Gracias a esto, demostraremos los resutados principales que nos permiten relacionar el tipo de homotopía de una variedad con el de un CW-complejo con tanta celdas de dimensión $\lambda$ como puntos críticos de índice $\lambda$ tenga una función de Morse sobre ella.

En un último capítulo enunciamos y demostramos las desigualdades de Morse, con las que podemos determinar importante información topológica de una varieda. En particular, dan cotas para los números de Betti de una variedad y una manera sencilla de computar su característica de Euler-Poincaré en base al número de puntos críticos de cada índice. Incluimos para este fin un capítulo previo de introduccion a la teoría de homología. En él construimos la homología simplicial de un espacio topológico triangulado, la homología celular para CW-complejos y la homología relativa.

%%%%%%%%%%%%%%%%%%%%%%%%%%%%%%%%
\fi


La teoría de Morse permite extraer de una variedad información sobre sus invariantes topológicos mediante el estudio de los puntos críticos de funciones diferenciables con valores reales.

Dada una función diferenciable $f:M\to \mathbb{R}$ sobre una variedad $M$, el lema de Morse establece que alrededor de todo punto crítico no degenerado existe un sistema local de coordenadas respecto del cual $f$ se expresa como una forma cuadrática. Esto tiene como consecuencia que los puntos críticos no degenerados son aislados y constituye el primer pilar de la teoría de Morse. Se dice que una función de este tipo es de Morse si todos sus puntos críticos son no degenerados. 

Si los conjuntos subnivel de una función $f$ de Morse son compactos, se puede relacionar el tipo de homotopía de la variedad con el de un CW-complejo con tantas celdas de dimensión $\lambda$ como puntos críticos de índice $\lambda$ tenga $f$. Por otro lado, las desigualdades de Morse determinan cotas para los números de Betti de una variedad y permiten computar su característica de Euler-Poincaré en función únicamente del número de puntos críticos de cada índice posible de una función de Morse.

En la presente memoria desarrollamos los fundamentos de esta teoría con el libro ``Morse Theory'' de John Milnor \cite{milnor1963} como fuente principal.